\begin{activity}{Numerical Demonstration: Isostatic Gravity Along a Cross Section}

\section*{Purpose}

In this demonstration, the instructor will use \texttt{isostasy\_gui.py} to build a lithospheric cross section along a chosen topographic profile, compute its gravity response under varying degrees of compensation and layered lithospheric structure, and compare the predicted gravity to observations. The purpose is to see single-column isostatic reasoning applied continuously along a real profile, and to meet the one genuinely new idea of the session: that compensation is not
all-or-nothing.

\section*{Learning Outcomes}

By the end of this demonstration, you should be able to:
\begin{itemize}
    \item Describe what the degree of compensation $C$ represents, and predict how changing it should affect the amplitude of the predicted gravity anomaly.
    \item Explain how the choice of layered lithospheric structure (crust only vs.\ crust + mantle lithosphere) affects the gravity prediction for the same topography.
    \item Use misfit between predicted and observed gravity to reason about how well (and where) a profile is locally compensated.
    \item Identify what a computational model needs beyond what was required for the idealized single-column calculations.
\end{itemize}

\subsection*{Background}

\paragraph{From columns to a profile.} Everything in this demonstration rests on the same isostatic mass balance used for a single column: at the depth of compensation, column-integrated mass must match between neighboring columns. The only difference here is scale --- instead of comparing two columns by hand, the program applies this balance (or a partial version of it, below) at every point along a chosen cross section, built from a loaded crustal thickness model rather than two fixed values.

\paragraph{Degree of compensation.} So far, every problem has assumed \emph{full} local compensation: the root exactly satisfies the isostatic mass balance. In reality, compensation is often partial, since lithospheric strength allows some of a load to be supported elastically rather than buoyantly. We describe this with a single parameter $C$, where
\[
    r(x) = C\, r_{\mathrm{Airy}}(x), \qquad 0 \le C \le 1,
\]
and $r_{\mathrm{Airy}}(x)$ is the fully compensated root thickness given by the isostatic (Airy) balance. $C=1$ recovers full local compensation; $C=0$ means the load is entirely supported elastically, with no compensating root at all.

\paragraph{Gravity of a real profile.} The infinite slab formula used for a single column assumes topography with no lateral edges. A real profile does not, so the program instead sums the gravity effect of many finite blocks along the cross section --- the same discrete 2-D prism approach introduced in Week~3 for irregularly shaped bodies, not a new formula.

\subsection*{App Overview}

\begin{itemize}
    \item Select two points to define a cross section across the loaded topography.
    \item Load a crustal thickness model for the same profile.
    \item Specify a layered lithospheric model (layer densities and thicknesses), following the same layered-column structure used for single-column isostatic balance.
    \item Set the degree of compensation $C$ and compute the predicted gravity anomaly along the profile.
    \item Compare the predicted gravity to observed gravity along the same profile, if loaded.
\end{itemize}

\section*{Tasks}

Follow along as the instructor builds a cross section and adjusts each parameter in turn. Answer the questions below as you go --- they are written to apply to whichever profile is selected.

\begin{questions}
    \question{Full compensation baseline}
    With $C=1$, does the predicted gravity anomaly broadly track or oppose the topography along the profile? Explain why, referring back to what perfect compensation implies for the combined gravity effect of topography and its compensating root.
    \answerbox[3cm]

    \question{Varying the degree of compensation}
    As the instructor decreases $C$ from 1 toward 0, predict first, then observe: does the amplitude of the predicted gravity anomaly increase or decrease? Explain in terms of how much compensating root mass remains at each value of $C$.
    \answerbox[3cm]

    \question{Finding the best-fit compensation}
    If observed gravity is available for this profile, adjust $C$ to find the value that best matches the predicted and observed gravity. What does that value of $C$ tell you about how this section of lithosphere is actually being supported?
    \answerbox[3cm]

    \question{Model dependence}
    With $C$ held fixed, the instructor will change the assumed crustal density and/or reference lithospheric thickness. How does the predicted gravity change? Relate this to the idea that the isostatic correction is model-dependent, not purely data-driven.
    \answerbox[3cm]

    \question{Adding a layer}
    The instructor will now add or modify a mantle lithosphere layer, as in the two-layer single-column model. Does the degree of compensation needed to fit the observed gravity change once this layer is included? Why might a two-layer model change your interpretation compared to a crust-only
    model?
    \answerbox[3cm]

    \question{Uniform compensation?}
    Does a single value of $C$ fit the observed gravity equally well across the whole profile, or does the fit degrade in some places (e.g.\ near a sharp topographic feature or a continent--ocean boundary)? What would that tell you about assuming one compensation state for an entire
    region?
    \answerbox[3cm]

    \question{What changed, computationally}
    Compared to a single-column, by-hand isostatic calculation, what additional ingredients did this model need to predict gravity along a real profile?
    \answerbox[3cm]

\end{questions}

\section*{Key Ideas}

\textbf{Compensation is a spectrum, not a switch.} Local isostatic balance and fully elastic support are two ends of a continuum; most real topography is compensated to some intermediate degree, and that degree can be estimated by comparing predicted and observed gravity.

\textbf{A single column-balance principle scales up without changing.} The same mass balance used for two columns by hand applies, unchanged, at every point along a continuous profile --- modelling a whole cross section is a matter of scale and bookkeeping, not new physics.

\textbf{Fitting a model is itself an inference about the Earth.} The value of $C$ (and the layered structure) that best matches observed gravity is not just a curve-fitting exercise --- it is an estimate of how a real piece of lithosphere is mechanically supported, with the same model-dependence caveats that apply to any isostatic correction.

\end{activity}